\section*{Devoir maison }
Un itinéraire est noté sous la forme d’une suite de lettres :
\begin{itemize}
    \item D pour un déplacement vers la droite d’une case.
    \item B pour un déplacement vers le bas d’une case.
\end{itemize}
Le nombre de caractères D est la longueur de l’itinéraire. Le nombre de caractères B
est sa largeur.
Ainsi l’itinéraire 'DDBDBBDDDDB' a pour longueur 7 et pour largeur 4. Sa représentation
graphique est :
\begin{verbatim}
S * *
    * *
      *
      * * * * *
               E
\end{verbatim}
avec :
\begin{itemize}
    \item S représente la case de départ (start). Ses coordonnées sont (0 ; 0)
    \item * représente les cases visitées
    \item E représente la case d’arrivée (end).
\end{itemize}

\begin{lstlisting}[language=Python]
class Chemin:
    def __init__(self, itineraire):
        self.itineraire = itineraire
        longueur, largeur = 0, 0
        for direction in self.itineraire:
            if direction == "D":
                longueur = longueur + 1
            if direction == "B":
                largeur = largeur +1
        self.longueur = longueur
        self.largeur = largeur
        self.grille = [['.' for i in range(longueur+1)]
                            for j in range(largeur+1)]

    def remplir_grille(self):
        i, j = 0, 0 # Position initiale
        self.grille[0][0] = 'S' # Case de départ marquée d'un S
        for direction in ...:
            if direction == 'D':
                ... = ... # Déplacement vers la droite
            elif direction == 'B':
                ... = ... # Déplacement vers le bas
            self.grille[i][j] = '*' # Marquer le chemin avec '*'
        self.grille[self.largeur][self.longueur] = 'E' # Case E

    def get_dimensions(self):
        ................................................
\end{lstlisting}

\begin{Exercise}[title={POO - Itinéraire}]
    \begin{lstlisting}[language=Python]
chemin_1 = Chemin("DDBDBBDDDDB") # on crée une instance de Chemin
a = chemin_1.largueur
b = chemin_1.longueur
    \end{lstlisting}
    \Question On exécute le code ci-dessous dans la console Python, préciser les valeurs contenues dans chacune des variables a et b.\\
    \vfill
    \emptybox{1cm}
    \Question Compléter la méthode remplir\_grille.
    \Question Écrire une méthode tracer\_chemin de la classe Chemin qui affiche une
représentation graphique d’un itinéraire (avec les étoiles comme dans l'énoncé).\\
    \vfill
    \emptybox{5cm}
    \Question Écrire une méthode get\_dimensions de la classe Chemin qui renvoie la
longueur et la largeur de l’itinéraire sous la forme d’un tuple.\\
    \vfill
    \emptybox{5cm}


\begin{lstlisting}[language=Python]
from random import choice

def itineraire_aleatoire(taille):
    itineraire = ''
    i = 0
    while .........................:
        direction = choice(["D", "B"])
        itineraire += .............
        ...............
    return itineraire
\end{lstlisting}
    \Question Compléter la fonction itineraire\_aleatoire qui crée un itinéraire aléatoire de taille N (N = longueur + largeur).\\
    \Question Combien y a-t-il de chemins possibles de taille N ? Justifiez.\\
    \vfill
    \emptybox{2cm}
\end{Exercise}